\chapter*{Conclusions}
\addcontentsline{toc}{chapter}{Conclusions}
\markboth{Conclusions}{Conclusions}


In this thesis we studied the impact of the uncertainties on the DM density profile of the Milky Way on the indirect detection constraints obtained from the Galactic Center.
We mainly investigated this region since the DM density is expected to rise towards the center, leading to a flux of gamma rays that is orders of magnitude larger than that predicted 
from the brightest dwarf galaxies. 
However, the GC is severely affected by the uncertainties on the density
profile and by the astrophysical background.

We first compared a set of representative density profiles, namely the Einasto, Burkert and cored NFW profiles, fixing their parameters through observational constraints and computing their $J$-factors (\cref{GC}).
The latter, integrated over a small region around the Galactic Center, span more than two orders of magnitude between the most peaked profile (Einasto) and the most cored one (Burkert), as summarized in \cref{tab:gc_jfactors}.
This has severe consequences on the upper bounds on the annihilation cross section. In particular, the upper bound derived under
the Burkert assumption is roughly a factor $\sim 4\times10^{2}$ weaker than the one obtained under Einasto.
We then considered a simple thermal relic model, in which DM annihilates into a pair of mediators that subsequently decay into $W$ bosons, with the annihilation cross section fixed by the relic abundance to $\langle\sigma v\rangle = 2.2\times10^{-26}~\mathrm{cm^3/s}$ for self-conjugate DM (\cref{sec:relic_constraints}).
Comparing this prediction with the H.E.S.S. Inner Galaxy Survey limits~\cite{Abdalla_2022}, provided for the NFW profile and rescaled to our halo parameters, we found that the model is not excluded at any DM mass, neither for the NFW profile nor for the cored one, whose limit was
obtained through a further rescaling of the $J$-factors (\cref{Cross_section}).

Subsequently, we studied the adiabatic growth of the supermassive black hole Sgr~A* at the center of the Milky Way (\cref{adiabatic_spike}):
the black hole redistributes the inner DM into a dense cusp, the spike, strongly enhancing the expected annihilation signal (\cref{Spikeprofiles}).
We introduced two benchmark spike profiles, the steep Gondolo--Silk (GS) spike (\cref{sec:GSspike}) and a milder one accounting for stellar heating (LSH, \cref{less_stellar_heating}), and compared them with our thermal relic model (\cref{sec:benchmarks}).
The spike enhances the $J$-factor of the innermost region by several orders of magnitude (\cref{tab:spike_jfactors}), and correspondingly strengthens the cross section limits derived from the H.E.S.S. observations of Sgr~A* (\cref{sec:spike_limits}).
For the GS benchmark, thermal relic DM is excluded over the whole mass range probed, $0.3 \lesssim m_{\mathrm{DM}}/\mathrm{TeV} \lesssim 100$, while for the LSH benchmark it is excluded only in the window $15 \lesssim m_{\mathrm{DM}}/\mathrm{TeV} \lesssim 69$ (\cref{Sigmav_spike}).


These results clearly show that the annihilation cross section constraints 
depend critically on both the poorly known inner DM distribution and the assumed spike benchmark.
Even within the NFW assumption, the choice of the spike benchmark alone moves the
excluded region from the entire mass range to a narrow window at high masses.
Nevertheless, it should be stressed that the two sources of uncertainty 
are not independent, since the spike develops from the inner halo 
and therefore inherits the uncertainties in both its inner slope and density normalization.